\documentclass[11pt]{article}
\usepackage{iftex}
\ifPDFTeX
  \usepackage[T1]{fontenc}
  \usepackage{lmodern}
\fi
\usepackage{amsmath,amssymb,amsthm}
\usepackage[margin=1in]{geometry}
\usepackage{booktabs,longtable,array}
\usepackage[numbers,sort&compress]{natbib}
\usepackage{microtype}
\usepackage{xurl}
\usepackage{hyperref}
\hypersetup{colorlinks=true,linkcolor=blue,citecolor=blue,urlcolor=blue,
  pdftitle={Sums of two rough numbers with prescribed Liouville signs},
  pdfauthor={Jizhou Guo}}
\Urlmuskip=0mu plus 1mu
\setlength{\emergencystretch}{2em}
\newtheorem{theorem}{Theorem}[section]
\newtheorem{lemma}[theorem]{Lemma}
\newtheorem{proposition}[theorem]{Proposition}
\newtheorem{corollary}[theorem]{Corollary}
\newtheorem{inputthm}[theorem]{Input}
\theoremstyle{definition}
\newtheorem{definition}[theorem]{Definition}
\newcommand{\e}{\mathrm e}
\newcommand{\E}{\mathbb E}
\newcommand{\Prob}{\mathbb P}
\newcommand{\Z}{\mathbb Z}
\newcommand{\N}{\mathbb N}
\newcommand{\C}{\mathbb C}
\newcommand{\R}{\mathbb R}
\newcommand{\ind}{\mathbf 1}
\newcommand{\Dpret}{\mathbb D}
\newcommand{\ph}{\varphi}
\newcommand{\leanname}[1]{\nolinkurl{#1}}
\newcommand{\norm}[1]{\left\lVert#1\right\rVert}
\newcommand{\abs}[1]{\left\lvert#1\right\rvert}
\DeclareMathOperator{\rad}{rad}
\DeclareMathOperator{\tr}{tr}
\DeclareMathOperator{\HS}{HS}
\DeclareMathOperator{\Rea}{Re}
\DeclareMathOperator{\cond}{cond}
\title{Sums of two rough numbers with prescribed Liouville signs}
\author{Jizhou Guo\\
  \href{mailto:mitsuha2021b@gmail.com}{\texttt{mitsuha2021b@gmail.com}}}
\date{October 2026}

\DeclareMathOperator{\supp}{supp}
\newcommand{\Pset}{\mathcal P}
\newcommand{\Sset}{\mathcal S}
\newcommand{\Rz}{R_z}
\newcommand{\wt}{W_{\mathbf t}}
\theoremstyle{remark}
\newtheorem{remark}[theorem]{Remark}
\begin{document}
\maketitle
\begin{abstract}
Let $\lambda$ be the Liouville function. We show that for every sufficiently large even integer $N$ and
each pair of signs $(\sigma,\tau)$ there are at least $cN/(\log N)^{1/4}$ integers $n\in[N/3,2N/3]$
such that $n$ and $N-n$ have no prime factor below $\exp((\log N)^{1/10})$, $\lambda(n)=\sigma$ and
$\lambda(N-n)=\tau$; moreover $n$ and $N-n$ may be taken squarefree with at most
$\lceil 2\log\log N\rceil$ prime factors. The proof combines the zero-free half-plane
$\operatorname{Re}s>7/8$ for Dirichlet $L$-functions, proved by OpenAI, with the weighted dilation
graphs of a second OpenAI paper. The graph estimates are needed for displacements that are multiples of
$N$ and for prime groups with the prime factors of $N$ removed, and we show that the proofs given
there apply to both.
\end{abstract}
\noindent\textit{2020 Mathematics Subject Classification.} Primary 11P32; Secondary 11N36, 11N25.

\section{Introduction}\label{sec:intro}
All logarithms are natural. For a positive integer $n$ put
$\lambda(n)=(-1)^{\Omega(n)}$, where $\Omega$ counts prime factors with
multiplicity. Write $P^-(n)$ for its least prime factor, with
$P^-(1)=\infty$, and $P^+(n)$ for its largest prime factor, with
$P^+(1)=1$. A number is $z$-rough if $P^-(n)>z$.

\begin{theorem}\label{thm:main}
There are absolute constants $c>0$ and $N_0$ such that, for every even
integer $N\ge N_0$ and every $(\sigma,\tau)\in\{-1,+1\}^2$,
\begin{equation}\label{eq:main}
\begin{split}
\#\{n:\ &N/3\le n\le 2N/3,\quad P^-(n)>z,\quad P^-(N-n)>z,\\
       &\lambda(n)=\sigma,\quad\lambda(N-n)=\tau\}
       \ge cN/(\log N)^{1/4},\qquad z=\exp((\log N)^{1/10}).
\end{split}
\end{equation}
\end{theorem}

\begin{theorem}\label{thm:postselection}
The representations in Theorem~\ref{thm:main} may also be required to have
$n$ and $N-n$ squarefree and
\begin{equation}\label{eq:factor-cutoff}
\Omega(n),\Omega(N-n)\le K(N):=\lceil2\log\log N\rceil.
\end{equation}
The lower bound is $cN/(\log N)^{1/4}$ with a possibly smaller absolute
constant $c>0$.
\end{theorem}

For even $N$ the number of $n\in[N/3,2N/3]$ satisfying just the two
roughness conditions has order
\begin{equation}\label{eq:rough-scale}
\mathfrak S(N)\frac{N}{(\log z)^2}
=\mathfrak S(N)\frac{N}{(\log N)^{1/5}},
\qquad
\mathfrak S(N)=2\prod_{p>2}\left(1-\frac1{(p-1)^2}\right)
                 \prod_{\substack{p\mid N\\p>2}}\frac{p-1}{p-2}.
\end{equation}
For completeness, the finite-residue argument in
Section~\ref{sec:weight}, applied just to the roughness product, gives its
Euler product with an error smaller than any fixed logarithmic power.
Mertens' theorem \cite[Section~2.2]{MV} identifies that product with
$\mathfrak S(N)(\log z)^{-2}$ up to absolute constant factors.
The factors of $N$ above $z$ change this comparison by
$1+O(L/(z\log z))$, where $L=\log N$.
The positive weight used below has mass at least
$cN/(\log N)^{59/250}$; we state the theorems with the exponent $1/4$.

In the author's earlier paper \cite{GuoEarlier}, all four sign patterns
of $(\lambda(n),\lambda(N-n))$ occur
$N/4+O(N/(\log N)^c)$ times for every sufficiently large $N$,
without a roughness restriction. Chen-type representations bound the
number of prime factors in the summands without prescribing both
parities. For example, Li and Liu \cite{LiLiu} prove that every sufficiently
large even $N$ is $p+rq$, where $p,q$ are primes, $r$ is $1$ or a prime,
and $r\le q^{0.9}$; the second summand can have either one or two
prime factors. Without roughness conditions, Mangerel
\cite{MangerelGRH} proves under GRH the $(-1,-1)$ pattern for every
sufficiently large even total and quantitative lower bounds for all
four patterns at sufficiently large prime totals.
CaptainSude \cite{CaptainSude} proves the $(-1,-1)$ pattern for every
even $N>2$. Meng \cite[Theorem~A]{Meng2026} proves that all four
patterns occur for every $N\ge2$ outside $\{2,3,4,5,6,9,10\}$, and
Pouly \cite[Theorem~1.1]{Pouly2026} proves that, for every fixed
$\varepsilon>0$, each pattern occurs at least
$c_\varepsilon N^{1-\varepsilon}$ times for every sufficiently large $N$.
Krishnamoorthy \cite{Krishnamoorthy2026} gives cancellation
for the binary Liouville convolution outside a small set of totals
and arbitrary logarithmic savings for its fixed higher convolutions.
The sign-pattern results of
\cite{MangerelGRH,CaptainSude,Meng2026,Pouly2026,Krishnamoorthy2026} impose no roughness
condition on the summands. A bound for sifted versions of the reflected two-point sum,
$\sum_{n\le N,\ P^-(n(N-n))>z_1}\lambda(n(N-n))$ in arithmetic progressions with $z_1$ a fixed power
of $N$, is stated as a hypothesis in \cite[Hypothesis~3.3]{Deligiannis2026}, and a sifted Liouville
cancellation at a fixed total is formulated in \cite{Arneth2026} as the remaining analytic problem of a
programme on the binary Goldbach problem. The bound $K(N)$ here grows with $N$.

We use the following inputs. The zero-free assertion is
Input~\ref{input:zf}; it is \cite[Theorem~1.1]{OpenAIQ78}, the theorem
of family~003 of the repository \url{https://github.com/openai/math},
and it is used with exactly that scope. Write [S] for
\cite{WeightedGraphs} and [E] for \cite{EvenParity}; both are taken
from the same repository at commit
\begin{center}\small
\nolinkurl{adc7f1241b42e322a6451854ab7e4b4c146bf78a}.
\end{center}
From [S] we use Definitions~3.1--3.3; Lemmas~3.4 and~3.6 together
with their proofs; Theorem~3.5 and the residue-model moment estimate
established in its proof (Sections~3.3--3.8, including
Lemmas~3.7--3.10); Corollary~3.11 and its proof; Theorem~4.1 and its
proof, including Lemma~4.3 and the Siegel--Walfisz estimate
[S, Theorem~2.2] used for (4.9); and Lemma~4.2 with (4.6).
Section~\ref{sec:graphs} states the graph inputs and their
order of choices. Proposition~\ref{prop:dilation} for dilated
displacements and Proposition~\ref{prop:punctures} for punctured prime
groups are obtained from the \emph{proofs} in [S], rather than its
statements. Their proofs verify, respectively, every use of the
displacement in Section~3 and every use of interval prime groups in
Section~4. Section~\ref{sec:cancellation} carries out, for the
reflected operator, the physical lift and shared-label normalization
of [S, Section~5.2] and the comparison reduction of [S, Section~5.3,
pp.~38--39], using the sign cancellation modeled on [E, Section~4].
We do not use the arithmetic endpoint theorem [S, Theorem~5.1].

The method retains the shared prime labels in a weighted graph and
then removes them using complete multiplicativity:
$\lambda(D_Rw)\lambda(D_Rw')=\lambda(w)\lambda(w')$.
Reflection is the identity $B_N=A_{-N}J_0$.
The two independent free products in each comparison term produce
rational frequency grids with denominators $Nr$. Character
orthogonality retains the factor $1/Q$ at $Q=Nr$; the zero-free input
bounds the short nonprincipal prime sum, and a small-prime modification
bounds the principal residual mean. The residual length exceeds $Q$
in the comparisons. At the raw scale $H'=1$ this inequality fails, so
the grid estimate is used only for comparison terms.

The exponent $1/10$ is not the limit of these estimates. The rough
cutoff must lie below the group primes, and the group exponents satisfy
$d<0.47$ in [S]. This bound keeps the residue moduli of moment length
$L^{0.52}$ equal to $\exp(o(L))$, and it prevents the argument from
reaching $z=N^\delta$ for fixed $\delta>0$.
The threshold $N_0$ is not effective with the Siegel--Walfisz input
used in [S, Section~4]. We make no assertion for a fixed number of
prime factors or for prime summands.

Section~\ref{sec:graphs} gives the graph inputs, dilation and reflection.
Section~\ref{sec:characters} establishes the character estimates.
Section~\ref{sec:grids} proves the grid energy estimate, including
a partial-summation verification of the smooth weights.
Section~\ref{sec:cancellation} proves reflected cancellation.
Section~\ref{sec:weight} constructs the positive weight and its mass.
Section~\ref{sec:theorems} deduces the two theorems.

\section{Graph inputs, dilation and reflection}\label{sec:graphs}
We use $L=\log x$ and $T=\log L$ in this section.

\begin{definition}[Prime groups and patterns, {\cite[Definitions~3.1--3.3]{WeightedGraphs}}]
\label{def:graphs}
Fix the data
\begin{equation}\label{eq:fixed-data}
0<a<b<c<d<0.47,\quad c_0,C_0,v_-,v_+>0,\quad
0<q_0<1,\quad \ell\in\N_{>0}.
\end{equation}
Let $\Pset_g$ be finite pairwise disjoint prime sets indexed by disjoint
small and big families $\mathcal G_{\rm small},\mathcal G_{\rm big}$.
Each family has between $c_0T$ and $C_0T$ groups. Put
\begin{equation}\label{eq:groups}
\begin{gathered}
\exp(L^a)\le p\le\exp(L^b)\quad(g\in\mathcal G_{\rm small}),\qquad
\exp(L^c)\le p\le\exp(L^d)\quad(g\in\mathcal G_{\rm big}),\\
v_-\le V_g:=\sum_{p\in\Pset_g}\frac1p\le v_+,\qquad
\mu_g(p)=\frac1{pV_g},\qquad
\omega_g(n)=\sum_{p\in\Pset_g}\ind_{p\mid n},\\
\Pset=\bigcup_g\Pset_g,\qquad s=|\mathcal G_{\rm small}|+
|\mathcal G_{\rm big}|,\qquad\omega_{\Pset}(n)=\sum_g\omega_g(n).
\end{gathered}
\end{equation}
Divisibility at every $n\in\Z$, including $0$, has its usual meaning.
For a fixed positive integer $J$ set $M=J+\ell$.
A pattern $\nu$ chooses $C_g\subset\{1,\ldots,J\}$, empty in small
groups and of size at most a fixed $m$ in big groups. Put
$t_g=\ell+|C_g|$. The slots in $\{1,\ldots,J\}\setminus C_g$
are shared between the ordered source and target lists.
The other $t_g$ endpoint slots are unshared. Independent auxiliary
labels with laws $\mu_g$ fill the $C_g$ slots in the edge product $D$,
which uses $J$ prime factors in each group.
Write $D=D_{\rm small}D_{\rm big}$.
The coefficients depend only on ordered big source, big target and
auxiliary big labels, and satisfy
\begin{equation}\label{eq:pattern-cost}
\sum_\nu\sup|K_\nu|\le L^{C_1},
\end{equation}
where $m,C_1$ are independent of $J$. Every operator is symmetrized
on both sides by averaging the $M!$ slot permutations independently
in each group.
\end{definition}

\begin{definition}[Ideal and physical operators]\label{def:operators}
The ideal Hilbert space is the product probability space of $M$
big-group labels with law $\mu_g$ in every big group; numerical
repetitions are allowed. A row retains shared labels, independently
draws unshared target and auxiliary free labels, and multiplies by
\begin{equation}\label{eq:ideal-multiplier}
K_\nu D_{\rm big}^{i\zeta}\e(\Theta D_{\rm big}),
\qquad \Theta,\zeta\in\R,\qquad \e(t)=\exp(2\pi it).
\end{equation}
Sum and symmetrize to obtain $T_\zeta(\Theta)$.
A physical state is $(n,\mathbf p)$, with $n\in\Z$ and $M$ ordered
distinct labels from each group, all dividing $n$. Its mass is
$\prod_g V_g^{-M}$, with counting measure in $n$.
For an integer $0<|k|\le L^{C_2}$ the row moves to $n+kD$,
sums compatible target lists with coefficient $\prod_gV_g^{-t_g}$,
and draws free labels with their $\mu_g$ laws. Each auxiliary free
label is forbidden from both endpoint lists, although free labels may
repeat one another. Its multiplier is
\begin{equation}\label{eq:physical-multiplier}
K_\nu D^{i\zeta}
q_0^{(\omega_{\Pset}(n)-Ms)/2}
q_0^{(\omega_{\Pset}(n+kD)-Ms)/2}.
\end{equation}
Call the symmetrized sum $A$.
Its adjoint reverses the displacement and interchanges the endpoint
labels, conjugating the multiplier. The joint list normalization is
$V_g^{-M-t_g}$ in either direction.
The vector $u_n$ equals $1$ on all lists at position $n$ and $0$
elsewhere; it is not normalized.
\end{definition}

\begin{inputthm}[Absolute bounds and transference, {\cite[Lemma~3.4,
Theorem~3.5, Corollary~3.11]{WeightedGraphs}}]\label{input:gt}
The absolute row and column sums of $A$ are bounded by $L^{C_{\rm abs}}$,
with $C_{\rm abs}$ independent of $J$ (it may depend on $m$, $C_1$ and the common fixed data); this also holds on the
independent-residue position space.
For each fixed $E>0$ there is $E_{\rm id}>0$, depending only on $E$
and the common fixed data, such that
\begin{equation}\label{eq:ideal-hypothesis}
\sup_\Theta\|T_\zeta(\Theta)\|\le L^{-E_{\rm id}}
\end{equation}
implies, for every sufficiently large fixed $J$, with
$R=\lceil L^{0.52}\rceil$, every fixed $\gamma>0$ and every positive dyad $I=[Y,2Y)$,
$x^\gamma\le Y\le x^{1/\gamma}$,
\begin{equation}\label{eq:source-moment}
\frac1{|I\cap\Z|}\sum_{n\in I\cap\Z}
\langle u_n,(AA^*)^Ru_n\rangle\le L^{-2RE}.
\end{equation}
The lower bound on $J$ may depend on $m,C_1$; the threshold in $x$
may depend on all fixed parameters, including $J,C_2,\gamma$.
Subject to these data and \eqref{eq:ideal-hypothesis}, the estimate
is uniform in groups, patterns, $k$ and $\zeta$.
There is no further bound on $\zeta$ when the ideal hypothesis holds.
If $f$ is supported on $I$, is independent of its list of marks,
$\sup|f_n|\le\exp(C_f\sqrt L)$, and
$\|f\|,\|g\|\le\sqrt Y L^{C_3}$, then
\begin{equation}\label{eq:source-pairing}
|\langle f,Ag\rangle|\ll YL^{-E+2C_3}.
\end{equation}
Only the first endpoint must be independent of its mark list.
\end{inputthm}

\begin{inputthm}[Residual ideal kernels, {\cite[Theorem~4.1, with (4.7)--(4.8) and the last two paragraphs of its proof, p.~34]{WeightedGraphs}}]
\label{input:ik}
Suppose initially that each big group is the set of all primes in an
interval in its permitted range. For each fixed $E_{\rm id}>0$ and
$C_4\ge0$ there are fixed $m,J_{\min},C_1$ and kernel parameters, independent
of $J$, such that every fixed $J\ge J_{\min}$ admits one raw pattern
$C_g=\varnothing,K_0=1$ and signed comparison patterns satisfying
\begin{equation}\label{eq:ik-bound}
\sup_{\Theta\in\R,\ |\zeta|\le L^{C_4}}
\|T_\zeta(\Theta)\|\le L^{-E_{\rm id}}.
\end{equation}
The family and coefficients are independent of the particular
$\Theta,\zeta$. Split the big groups into two blocks and designate $m$
of the first $J$ slots in each big group as probes. Each comparison frees
nonempty probe sets $\mathcal A,\mathcal B$ in the two blocks, with
coefficient
\begin{equation}\label{eq:ik-coefficient}
-(-1)^{|\mathcal A|+|\mathcal B|}
 K_{\mathcal A}(D_{\mathcal A},X_{\mathcal A})
 K_{\mathcal B}(D_{\mathcal B},X_{\mathcal B}).
\end{equation}
Here $D_{\mathcal A},D_{\mathcal B}$ are independent free-label
products; $X_{\mathcal A}$ uses target labels and $X_{\mathcal B}$
source labels. No shared label or final $\ell$ label enters these
coefficients. For a slot set $\Lambda$,
\begin{equation}\label{eq:cell-kernel}
K_\Lambda(D,X)=
\sum_{j:\nu_j\ge\xi}
\frac{\ind_{\log D\in I_j}\ind_{\log X\in I_j}}{\nu_j}
\sum_{\substack{\chi\ {\rm primitive}\\\cond(\chi)\le Q_0}}
\overline{\chi(D)}\chi(X),
\end{equation}
where $I_j=[j\mathfrak h,(j+1)\mathfrak h)$ has width $\mathfrak h\le1$,
$\nu_j$ is the product-law cell probability, and
$\mathfrak h^{-1},\xi^{-1},Q_0$ are fixed powers of $L$.
The conductor-$1$ character is included. The number and total coefficient
mass of the expanded terms are fixed powers of $L$, independent of $J$.
Every free product contains a big prime.
\end{inputthm}

The order in these inputs is substantive.
In [S, (3.37), p.~23] one fixes
$G>E+C_v+3$ and $E_{\rm id}>G+C_s+2$ before increasing $J$.
In [S, pp.~33--34] one fixes the probe count $m$, then the comparison
accuracy, then the kernel parameters in (4.20) and their coefficient
exponent $C_1$, and only then imposes a lower bound on $J$.
Thus $J$ may subsequently be increased to satisfy transference
without changing any kernel precision. We will keep this order.

\begin{proposition}[Coprime integer dilation]\label{prop:dilation}
Keep Definitions~\ref{def:graphs}--\ref{def:operators} and the fixed
bound $0<|k|\le L^{C_2}$. Let $h\ne0$ be any integer coprime to every
group prime, and replace every physical displacement $kD$ by $hkD$,
leaving the coefficients and all mark rules unchanged. Call this
operator $A_h$. Under \eqref{eq:ideal-hypothesis}, with the same order
of fixed parameters, one has, for every fixed $\gamma>0$ and every positive dyad $[Y,2Y)$ with $x^\gamma\le Y\le x^{1/\gamma}$,
\begin{equation}\label{eq:dilated-moment}
\sum_{Y\le n<2Y}\langle u_n,(A_hA_h^*)^Ru_n\rangle
\ll YL^{-2RE}
\end{equation}
and, under the endpoint hypotheses of Input~\ref{input:gt},
\begin{equation}\label{eq:dilated-pairing}
|\langle f,A_hg\rangle|\ll YL^{-E+2C_3}.
\end{equation}
The absolute row and column bounds of Input~\ref{input:gt} hold for
$A_h$ with the same $C_{\rm abs}$, also on the independent-residue
position space. The constants and threshold are independent of $h$
and its size. This proposition is deduced from the proofs of
Lemmas~3.4 and~3.6, the residue-model moment estimate in the proof of
Theorem~3.5, and the proof of Corollary~3.11 of [S].
\end{proposition}
\begin{proof}
We spell out the dependence on the displacement in [S, Section~3]; in this proof, Definition, Lemma, Theorem and Corollary numbers and equation numbers of the form (3.$n$) are those of [S].
Definition~3.3 and the adjoint computation (3.3)--(3.4), pp.~11--12,
replace $n+kD$ by $n+hkD$; reversal still keeps $D$ fixed and the
joint normalization unchanged. The absolute estimate (3.5) in
Lemma~3.4, p.~12, counts target marks and uses only divisibility and
translation of the position; its row and column bounds therefore hold
for $A_h$, also on the independent-residue position space.
The range of $k$ in Theorem~3.5, p.~12, is the restriction being extended;
its residue proof uses that restriction only in the ways checked below.

In the root-path expansion, (3.8), p.~13, the offsets become $hb_i$,
where $b_0=0$ and $b_{i+1}-b_i=\pm kD_i$. Return is
$hb_{2R}=0$, equivalent to $b_{2R}=0$. Thus the closed numerical paths
and scalar coefficients are exactly those for $h=1$.
Lemma~3.6 and (3.9)--(3.10), pp.~13--14, specify congruences
$n+hb_i\equiv0\pmod p$. Their active and Bonferroni moduli use
the same primes, have size $\exp(o(L))$, and each CRT count on
$[Y,2Y)$ has rounding error $O(1)$ independently of the representatives
$hb_i$. The closure phase (3.11), p.~14, can be kept as
$\e(\theta b_{2R})$, since the return indicators coincide.

In the independent-residue law the measure-preserving map
$n_p\mapsto h^{-1}n_p$ changes every predicate
$n_p+hb_i=0$ to $n_p+b_i=0$. It preserves active visits,
repeated residues and all damping. Hence every numerical-path
expectation, including the absolute sum that controls replacement
error, is exactly the $h=1$ expectation. Consequently the signed
residue-model path sum for $A_h$, with the return condition
$b_{2R}=0$, equals the one for $h=1$, which the proof of
[S, Theorem~3.5] bounds by $o(L^{-2RE})$ (Sections~3.3--3.8 and the
completion on p.~25). The remainder of this paragraph and the next
two paragraphs record that each step of that proof also applies
directly to the displacement $hkD$. The lifespan tests
(3.13)--(3.14), pp.~14--15, are preserved as well.
The assertion on p.~15 that every group prime exceeds $|k|$
is used to infer $p\mid kD\Rightarrow p\mid D$.
For the enlarged displacement, $p>|k|$ and $p\nmid h$ give
$p\mid hkD\Rightarrow p\mid D$. Its converse exclusion for
unshared entry and exit labels is therefore unchanged.

In the edge-memory operations on p.~16, the translation
$\Delta=\pm kD$ becomes $\Delta=\pm hkD$, while the phase remains
$\e(\pm\theta kD)$ when the closure phase is kept as
$\e(\theta b_{2R})$.
Translations remain measure preserving. The exact memory identity
(3.17), p.~17, uses only these congruences and the just-verified
active/shared implication. The reversed memory measures
(3.23)--(3.25), p.~19, remain unchanged under any translation $\Delta$;
their absolute bounds, (3.26)--(3.28), p.~20, have no size dependence.
For rare transfers, (3.31) and Lemma~3.9, pp.~20--21, the congruence
is $r+\sigma_i hkD_{\rm big}F/U=0\pmod p$.
The omitted small-label products satisfy
$U\le\exp(\ell|\mathcal G_{\rm small}|L^b)<\exp(L^c)\le p$;
they are distinct modulo $p$, and $h,k,F,D_{\rm big}$ are invertible
there whenever $p\nmid D$. At most one omission choice is allowed,
exactly as in the source.

The clean-edge Fourier reduction (3.36) and Lemma~3.10, pp.~21--22,
only change the real frequency $\Theta$ in $\e(\sigma_i\Theta
D_{\rm big})$. The ideal hypothesis is uniform in every real $\Theta$.
The equality-rank expansion (3.39)--(3.47), pp.~23--25, changes no
birth labels, point masses or equality phases; its unaffected edges
retain this same clean estimate. The completion on p.~25 thus gives
the same residue-model moment. Restoring the uniform residue error
with a slightly stronger fixed precision proves
\eqref{eq:dilated-moment}. This is a full-root CRT average, with no
assertion of residue independence along a short progression modulo $h$.

It remains to adapt the physical interval bound $H$ in [S, (3.49),
p.~26]. Choose
\begin{equation}\label{eq:quotient-radius}
H_0\ge\max|kD|,\qquad \log H_0=O_J(TL^d).
\end{equation}
For each residue class modulo $|h|$, partition its integer quotient
coordinate into blocks of $\lceil H_0\rceil$ consecutive positions.
Write ${\rm block}_j$ for the part of one such block in $[Y,2Y)$, let
$f_j$ be the restriction of $f$ to ${\rm block}_j$, and let $g_j$ be
the restriction of $g$ to the union of the full quotient block containing
${\rm block}_j$ and the two adjacent full quotient blocks of that same
residue class, without intersecting these blocks with $[Y,2Y)$. Every edge stays in the
class and moves at most $H_0$ quotient positions. Therefore
\begin{equation}\label{eq:quotient-blocks}
\langle f,A_hg\rangle=\sum_j\langle f_j,A_hg_j\rangle,
\qquad \sum_j\|g_j\|^2\le3\|g\|^2.
\end{equation}
The positive semidefinite matrix
$(\langle u_n,(A_hA_h^*)^Ru_m\rangle)_{n,m\in{\rm block}_j}$
has largest eigenvalue at most its trace $d_j$. Mark independence gives
\begin{equation}\label{eq:block-trace}
\langle f_j,(A_hA_h^*)^Rf_j\rangle
\le\left(\sum_{n\in{\rm block}_j}|f_n|^2\right)d_j
\le\lceil H_0\rceil\sup|f_n|^2d_j.
\end{equation}
Spectral H\"older and H\"older in $j$, with exponents
$2,2R/(R-1),2R$, give
\begin{equation}\label{eq:block-holder}
|\langle f,A_hg\rangle|
\le\sqrt3\,\|g\|\|f\|^{1-1/R}
(\lceil H_0\rceil\sup|f_n|^2)^{1/(2R)}
\left(\sum_jd_j\right)^{1/(2R)}.
\end{equation}
The middle factor is bounded, since $d<0.47<0.52$ and $0.5<0.52$.
The blocks partition the support dyad, so their traces sum to
\eqref{eq:dilated-moment}. The norm hypotheses prove
\eqref{eq:dilated-pairing}. Empty blocks, including the case
$|h|>2Y$, cause no change.
\end{proof}

\begin{proposition}[Reflection]\label{prop:reflection}
Let $(J_0g)(n,\mathbf p)=g(-n,\mathbf p)$. For $h\ne0$ coprime to
the group primes, let $B_h$ have the source mark transitions with
physical target $hkD-n$. Then
\begin{equation}\label{eq:reflection}
B_h=A_{-h}J_0,\qquad B_hB_h^*=A_{-h}A_{-h}^*.
\end{equation}
Consequently \eqref{eq:dilated-pairing} and the absolute row and column
bounds of Proposition~\ref{prop:dilation} hold with $B_h$ in place of
$A_h$.
\end{proposition}
\begin{proof}
Reflection preserves divisibility, state measure and degree damping,
so $J_0$ is a unitary permutation of the states and leaves absolute
row and column sums unchanged. The translation target $n-hkD$ becomes $hkD-n$
after applying $J_0$, with exactly the same admissible lists.
The squared-operator identity follows from $J_0J_0^*=1$.
The second endpoint in the translation interface may be supported
at negative integers, since the physical space is all of $\Z$.
We only use Liouville values at the positive endpoints before applying
$J_0$.
\end{proof}

\begin{proposition}[Punctured prime intervals]\label{prop:punctures}
Set $x=N$. Remove the primes dividing $N$ from each prime interval
used for a big group. If the resulting groups satisfy
\eqref{eq:groups}, the conclusion and parameter order of
Input~\ref{input:ik} hold for them.
The transference inputs already allow arbitrary finite groups.
\end{proposition}
\begin{proof}
For any fixed $\alpha>0$, deletion has reciprocal mass at most
\begin{equation}\label{eq:deleted-mass}
\sum_{\substack{p\mid N\\p\ge\exp(L^\alpha)}}\frac1p
\le\frac{L}{L^\alpha\exp(L^\alpha)}.
\end{equation}
Indeed at most $L/L^\alpha$ distinct prime factors can lie above this
threshold. This is smaller than every fixed negative power of $L$,
also after any further interval intersection. For big groups take
$\alpha=c$; $\alpha=a$ controls both kinds of groups.

Here is the full check of the interval uses in [S, Section~4]; in this proof, Theorem, Lemma and equation numbers of the form (4.$n$) are those of [S].
Its opening assumption on p.~27 says that intersecting a big group
with another interval is still an interval of primes. After deletion
it is that interval with only primes dividing $N$ missing.
For a selected slot set $\Lambda$, write $P_\Lambda$ for its label
product, $k_g$ for its number of slots in group $g$, and $a_p$ for
the multiplicity of $p$ in a supported product $n$.
Let $\rho_\Lambda(n)=\Prob(P_\Lambda=n)$, zero outside the support.
In (4.3), p.~28, unique factorization gives, on this support, the mass
\begin{equation}\label{eq:product-collapse}
\rho_\Lambda(n)
=\frac1n\prod_g\frac{k_g!}{V_g^{k_g}\prod_{p\in\Pset_g}a_p!}
\le \frac{L^\kappa}{n},\qquad
\kappa=C_0\{\log(m!)+m\log\max(1,v_-^{-1})\},
\end{equation}
for at most $m$ selected slots per big group. This uses no interval
distribution and is unchanged, including repeated labels.
The fiber Cauchy bound (4.4), p.~28, is
\begin{equation}\label{eq:fiber-collapse}
\sum_n|\E[f\ind_{P_\Lambda=n}]|^2
\le \frac{L^\kappa}{c_1U}\|f\|_2^2
\quad(c_1U\le P_\Lambda\le c_2U);
\end{equation}
it follows from the same mass bound for arbitrary tuple tests.

The log-cell definitions (4.7) and bounds (4.8), p.~29, become
\begin{equation}\label{eq:cell-bounds}
N_\Lambda\ll_m1+TL^d/\mathfrak h,\qquad
\sup|K_\Lambda|\le Q_0^2\xi^{-1},\qquad
\sup_X\E_D|K_\Lambda(D,X)|\le Q_0^2.
\end{equation}
Support size, retained cell mass and the cancellation of $\nu_j$
give these inequalities for the new product law.

The only distribution estimate is (4.9), pp.~29--30:
\begin{equation}\label{eq:cell-offdiagonal}
\left|\E\left[\ind_{\log D\in I_j}
                  \overline{\chi(D)}\chi'(D)\right]\right|
\ll_{D_*} L^{-D_*}\quad(\chi\ne\chi').
\end{equation}
Fix all but one selected label. The remaining prime lies in the
intersection of its original interval and a cell-imposed interval.
The character $\overline\chi\chi'$ is nonprincipal on a modulus
at most $Q_0^2$, by uniqueness of primitive induction.
The reciprocal prime sum used in [S, (4.9)] is valid on the full
intersection by [S, Theorem~2.2]; removing $p\mid N$ changes it
by at most \eqref{eq:deleted-mass}. Division by $V_g\ge v_-$
and averaging the other labels preserve arbitrary precision.
Every remaining prime exceeds $Q_0^2$, so all conditional diagonal
Gram entries in (4.10), p.~30, are exactly $1$.
Thus (4.11), p.~30, retains the bound
\begin{equation}\label{eq:gram-bound}
1+O_{D_*}(Q_0^2\xi^{-1}L^{-D_*})\le2.
\end{equation}
The character-reproduction formula (4.18), p.~32, remains
\begin{equation}\label{eq:character-reproduction}
\E_D K_\Lambda(D,X)\psi(D)
=\psi(X)+O_{D_*}(Q_0^2\xi^{-1}L^{-D_*})
\end{equation}
on retained cells: its diagonal term is exact and the others use
\eqref{eq:cell-offdiagonal}. Applying it twice with the rational
character expansion (4.15) preserves (4.19), p.~32, whose error is
\begin{equation}\label{eq:rational-reproduction}
O_{D_*}(Q_0^{5/2}\xi^{-1}L^{-D_*}).
\end{equation}

For the remaining displays, the integer bilinear estimate (4.5)--(4.6),
pp.~28--29, permits coefficients vanishing on any subset of containing
intervals. The cell support comparison and derivative error
(4.13)--(4.17), pp.~30--31, use only product sizes,
$\mathfrak h$, and the new cell probabilities; discarded cells still
have total mass at most $(N_{\mathcal A}+N_{\mathcal B})\xi$.
The parameter inequalities (4.20), p.~32, absorb
\eqref{eq:deleted-mass} at any fixed precision before $J$.
The symmetric-probe fractions and cancellations
(4.21)--(4.26), pp.~32--34, use a product probability law and
within-group permutation invariance, both of which persist.
These are all the interval dependencies: none asserts distribution
for an arbitrary subset of a prime interval. The proof of [S,
Theorem~4.1] therefore gives the proposition, uniformly in $N$.
\end{proof}

\section{Zero-free input and character estimates}\label{sec:characters}

\begin{inputthm}[Family 003, {\cite[Theorem~1.1]{OpenAIQ78}}]\label{input:zf}
For every positive integer $q$, every Dirichlet character $\chi\pmod q$
and every $s\in\C$ with $\Rea s>7/8$, one has $L(s,\chi)\ne0$,
apart from the principal character's pole at $s=1$.
\end{inputthm}
The Lean declaration is
\begin{center}\small
\path{OAI.DirichletCharacter.LFunction_ne_zero_of_seven_eighths_lt_re}.
\end{center}
The companion declaration
\path{OAI.riemannZeta_ne_zero_of_seven_eighths_lt_re}
gives the zeta assertion away from its pole.

\begin{lemma}[Liouville character power bound]\label{lem:power}
For fixed $C_1>0$ and every fixed $\alpha>7/8$,
\begin{equation}\label{eq:power-bound}
\sum_{n\le X}\lambda(n)\chi(n)n^{it}\ll_{C_1,\alpha}X^\alpha
\qquad(q,|t|\le X^{C_1}),
\end{equation}
uniformly in all characters modulo $q$, including principal and
imprimitive characters.
\end{lemma}
\begin{proof}
We give the analytic details, using the classical Borel--Carath\'eodory and three-lines inequalities
and the Dirichlet-series and contour tools of \cite[Chapters~5, 10--13]{MV}.
For primitive nonprincipal $\chi$, Input~\ref{input:zf} permits a
holomorphic branch $F(s)=\log L(s,\chi)$ on $\Rea s>7/8$,
agreeing with the Euler logarithm on $\Rea s>1$.
Partial summation of the character sum, bounded by $q$, gives
$|L(s,\chi)|\le q|s|/\Rea s$ for $\Rea s>0$.
Borel--Carath\'eodory on discs centered at $2+it$, staying a fixed
distance to the right of $7/8$, bounds
$|F(7/8+\delta/2+it)|\ll_\delta\log(q(|t|+3))$.
On $\Rea s=1+\delta/2$ the Euler logarithm is $O_\delta(1)$.
Apply three-lines to
$F(s)\exp((s-it_0)^2/\mathcal R^2)$, where
$\mathcal R=q+|t_0|+3$. The Gaussian gives finite boundary suprema.
At $s=\sigma+it_0$, $\sigma\ge7/8+\delta$, this yields
\begin{equation}\label{eq:logarithm-bound}
|F(s)|\ll_\delta
[\log(q(|t_0|+3))]^{1-\eta_\delta},
\qquad \eta_\delta>0,
\end{equation}
and hence, for every $\varepsilon>0$,
\begin{equation}\label{eq:reciprocal-growth}
L(s,\chi)^{-1}\ll_{\delta,\varepsilon}
(q(|t|+3))^\varepsilon\qquad(\Rea s\ge7/8+\delta).
\end{equation}
For the primitive principal character use the nonvanishing analytic
function $Z(s)=\zeta(s)(s-1)/(s+1)$, whose canceled value at $1$ is
nonzero. Its logarithm has the same bounds; the factor
$(s-1)/(s+1)$ in $1/\zeta(s)$ is bounded by $1$ on $\Rea s\ge0$.
For imprimitive characters, the inverse local-factor logarithm is
$O_\delta(\sum_{p\mid q}p^{-7/8-\delta})=o(\log q)$.
To see uniformity of the last assertion, fix a large prime threshold:
the finitely many smaller primes cost $O(1)$, and the contribution of
larger primes, divided by $\log q$, is arbitrarily small.
It is absorbed by another arbitrarily small conductor power.

The twisted series is
\begin{equation}\label{eq:liouville-series}
\sum_{n\ge1}\frac{\lambda(n)\chi(n)n^{it}}{n^s}
=\frac{L(2s-2it,\chi^2)}{L(s-it,\chi)}.
\end{equation}
The numerator is absolutely bounded on a fixed line right of $7/8$.
Truncated Perron at a half-integer cutoff and height $\mathcal T$
has error $O(X\log^2X/\mathcal T+1)$.
Move to $\Rea s=7/8+\delta<\alpha$; the vertical term is
\begin{equation}\label{eq:perron-vertical}
O\left(X^{7/8+\delta}
       (q(\mathcal T+|t|+3))^\varepsilon\log\mathcal T\right),
\end{equation}
and the horizontal integrals are smaller after a sufficiently large
fixed-power height is chosen. There is no principal pole in the ratio:
the pole of $L(s-it,\chi)$ becomes a zero of its inverse.
Choose the reciprocal-growth exponent $\varepsilon$ after $C_1$,
the height exponent and the margin $\alpha-7/8-\delta$.
This proves \eqref{eq:power-bound}. The same reasoning applies when
$q,|t|\le N^C$ and $N^\eta\le X\le N^{1/\eta}$ by enlarging $C_1$.
\end{proof}

\begin{lemma}[Small-prime modifications]\label{lem:modifications}
Fix $C>0$, $0<d<1$, and $L=\log N$.
Let $Z_0\le\exp(L^d)$ and let $f$ be real and completely
multiplicative, possibly depending on $N$, with
$f(p)\in[-1,0]$ for $p\le Z_0$ and $f(p)=-1$ otherwise.
For every fixed $\eta,A>0$,
\begin{equation}\label{eq:modified-mean}
\sup_{N^\eta\le X\le N^{1/\eta}}
\sup_{\substack{q\le N^C\\|t|\le N^C}}
\sup_{\chi\bmod q}
\left|\sum_{n\le X}f(n)\chi(n)n^{it}\right|
\ll_{C,d,\eta,A}XL^{-A}.
\end{equation}
This includes every character and real $t$; the same bound holds on
any subinterval of $[X,2X]$.
\end{lemma}
\begin{proof}
Write $f(p)=-u_p$, $0\le u_p\le1$. The exact convolution $f=\lambda*g$
has local coefficients
\begin{equation}\label{eq:modified-euler}
\sum_{j\ge0}g(p^j)v^j=\frac{1+v}{1+u_pv},\qquad
g(p^j)=(1-u_p)(-u_p)^{j-1}\quad(j\ge1)
\end{equation}
at modified primes, and $g(p^j)=0$ at all others. For $0<v<1$,
\begin{equation}\label{eq:modified-local-bound}
\sum_{j\ge1}|g(p^j)|v^j
=\frac{(1-u_p)v}{1-u_pv}\le v.
\end{equation}
Mertens' estimates give
$\sum_n|g(n)|/n\ll\log(2Z_0)$.
For large $Z_0$, take $\delta=1/\log Z_0$ and apply Rankin:
\begin{equation}\label{eq:modified-tail}
\sum_{n>\sqrt X}\frac{|g(n)|}{n}
\le X^{-\delta/2}\prod_{p\le Z_0}(1+p^{-1+\delta})
\ll X^{-1/(2\log Z_0)}(\log(2Z_0))^{\mathrm e}.
\end{equation}
The last estimate uses $p^\delta\le\mathrm e$ and the harmonic
prime form of Mertens' theorem.
For $n\le\sqrt X$, Lemma~\ref{lem:power} with $\alpha=15/16$ gives
\[
X^{15/16}\sum_{n\le\sqrt X}|g(n)|n^{-15/16}
\le X^{31/32}\sum_n|g(n)|/n.
\]
Thus the complete sum is
\begin{equation}\label{eq:modified-total}
\ll X^{31/32}\log(2Z_0)
+X\exp\left(-\frac{\log X}{2\log Z_0}\right)
                        (\log(2Z_0))^{\mathrm e}.
\end{equation}
Both terms are $O_A(XL^{-A})$, uniformly in $u_p,\chi,t$.
Bounded $Z_0$ uses a fixed Rankin exponent; no modifications gives
$f=\lambda$ directly. Differences of prefixes give subintervals.
\end{proof}

\begin{corollary}[Rough one-factor means]\label{cor:rough-mean}
Put $\Rz(n)=\ind_{P^-(n)>z}$. For fixed $C_0>0$,
$0<\rho_0<1$ and $H>0$, uniformly for
$2\le z\le\exp((\log X)^{\rho_0})$, $q,|t|\le X^{C_0}$ and all
characters modulo $q$,
\begin{equation}\label{eq:rough-mean}
\sum_{n\le X}\lambda(n)\Rz(n)\chi(n)n^{it}
\ll_{C_0,\rho_0,H}X(\log X)^{-H}.
\end{equation}
\end{corollary}
\begin{proof}
Let $\mu$ be the M\"obius function, equal to $0$ on nonsquarefree
integers and to $(-1)^{\omega(n)}$ on squarefree integers, where
$\omega(n)$ counts distinct prime factors.
The coefficient identity, including prime powers, is
\begin{equation}\label{eq:rough-convolution}
\lambda\Rz=\lambda*g_z,\qquad
g_z(d)=\mu^2(d)\ind_{P^+(d)\le z}.
\end{equation}
At $p\le z$ the factor $1+p^{-s}$ cancels the Liouville Euler
factor; at other primes it is $1$.
Writing $\alpha=15/16$ and taking the conductor exponent
$C_1=2C_0$ on the short convolution gives
\begin{equation}\label{eq:rough-small-part}
\sum_{d\le\sqrt X}g_z(d)(X/d)^\alpha
\le X^{(1+\alpha)/2}\sum_dg_z(d)/d
\ll X^{31/32}\log(2z).
\end{equation}
With $\varepsilon_z=1/\log z$ the other part satisfies
\begin{equation}\label{eq:rough-tail}
\sum_{d>\sqrt X}\frac{g_z(d)}d
\le X^{-1/(2\log z)}
 \prod_{p\le z}(1+p^{-1+1/\log z})
\ll X^{-1/(2\log z)}(\log(2z))^{\mathrm e}.
\end{equation}
The finite product is valid even if $\varepsilon_z>1$ for small
fixed $z$. The total bound is
\begin{equation}\label{eq:rough-total}
C\left\{X^{31/32}\log(2z)
+X\exp\left[-\frac{\log X}{2\log z}\right]
 (\log(2z))^{\mathrm e}\right\}.
\end{equation}
Since $\log z\le(\log X)^{\rho_0}$ with $\rho_0<1$ fixed, both terms are $O_H(X(\log X)^{-H})$, which proves \eqref{eq:rough-mean}.
\end{proof}

\begin{lemma}[Short prime character sums]\label{lem:prime}
For fixed $C>0$, $0<a<b<1$, and
$\exp(L^a)\le P\le\exp(L^b)$, every nonprincipal character modulo
$q\le N^C$, every real $|t|\le N^C$, and every interval
$I\subset[P,2P]$ satisfy
\begin{equation}\label{eq:prime-bound}
\left|\sum_{p\in I}\chi(p)p^{it}\right|
\ll_C P^{15/16}L^{C_1},
\end{equation}
where $C_1$ is fixed (with the fixed exponent ranges understood).
The same estimate holds after omitting primes dividing $N$.
For a principal character we use only $O(P)$.
\end{lemma}
\begin{proof}
Cauchy's formula for the analytic logarithm in
Lemma~\ref{lem:power}, or its standard zero partial-fraction expansion,
gives $-L'/L(s,\chi)=O(\log^2(q(|\Im s|+3)))$ on
$\Rea s=15/16$ and the horizontal segments, for nonprincipal
characters. Apply truncated Perron to
$-L'/L(s-it,\chi)$, initially with height $N^{C+10}$, increasing its
fixed exponent if required. Shift to $\Rea s=15/16$.
The pole at $1+it$ is absent precisely for nonprincipal $\chi$.
The zero-free input gives a fixed distance from all zeros.
The remaining line contributes $O(P^{15/16}L^{C_1})$.
Truncation and horizontal errors are polynomially small after the
height is enlarged, even for $P=\exp(L^a)$.
Proper prime powers cost $O(P^{1/2}L^{C_1})$.
Partial summation removes $\log p$; differences of prefixes treat $I$.
Imprimitive Euler factors and deletion of $p\mid N$ contribute at
most a fixed power of $L$, absorbed in \eqref{eq:prime-bound}.
Thus the modulus need only be polynomial in $N$, not in $P$.
A principal pole would supply a main term and is never discarded.
\end{proof}

\section{Fourier energy on rational grids}\label{sec:grids}
From now on $L=\log N$. Fix the separated ranges
\begin{equation}\label{eq:separated-ranges}
\begin{gathered}
z=\exp(L^{1/10}),\qquad
\Pset\subset[\exp(L^{0.12}),\exp(L^{0.18})]
       \cup[\exp(L^{0.35}),\exp(L^{0.40})],\\
\Sset=\{p:\exp(L^{0.22})<p\le\exp(L^{0.28}),\ p\nmid N\}.
\end{gathered}
\end{equation}
Assume $\Pset$ contains no prime factor of $N$ and
$\sum_{p\in\Pset}1/p=O(\log L)$. Write $n=n_{\Pset}n_*$,
where $n_{\Pset}$ is the full $\Pset$-part, including prime powers.
Use $\omega_{\Sset}$ for the number of distinct $\Sset$-prime
divisors and $\Omega_{\Sset}$ for their number with multiplicity.

\begin{lemma}[Grid energy]\label{lem:grid}
Let $c(g)$ be a complex function supported on $\Pset$-smooth integers,
that is, integers all of whose prime factors lie in $\Pset$,
and bounded by $L^C$. It may depend on $N$ and all fixed comparison
parameters. Define
\begin{equation}\label{eq:grid-sequence}
a(n)=\psi(n/Y)n^{iv}\lambda(n_*)\Rz(n_*)
 \ind_{(n,N)=1}\ind_{\omega_{\Sset}(n_*)>0}\,c(n_{\Pset}),
\end{equation}
where $\psi$ is supported on a fixed compact subinterval of
$(0,\infty)$,
$\|\psi^{(j)}\|_\infty=O_j(L^{C(j+1)})$, and $v\in\R$, $|v|\le L^C$.
Fix $C_0,C_2>0$. Uniformly for
\begin{equation}\label{eq:grid-scales}
Y=NH',\qquad
\exp(2L^{0.35})/4\le H'\le\exp(C_2L^{0.40}\log L),
\qquad r\in\N_{>0},\ r\le L^{C_0},\quad Q=Nr,
\end{equation}
and every $A>0$,
\begin{equation}\label{eq:grid-energy}
\int_{|\beta|\le L^{C_0}/Y}
 \sum_{b\bmod Q}|\widehat a(b/Q+\beta)|^2\,d\beta
\ll_{A,C,C_0,C_2} YL^{-A},
\qquad \widehat a(\theta)=\sum_na(n)\e(\theta n).
\end{equation}
Consequently the energy on the union of these arcs for all
$r\le L^{C_0}$ is $O_A(YL^{-A})$, after increasing intermediate precision.
\end{lemma}
\begin{proof}
We begin with a prime extraction that does not assume squarefreeness.
For $\omega_{\Sset}(n_*)>0$ one has the exact identity
\begin{equation}\label{eq:reserved-extraction}
-\int_0^1\sum_{\substack{p\in\Sset\\p\mid n_*}}
\lambda(n_*/p)\Rz(n_*/p)
 u^{\Omega_{\Sset}(n_*/p)}\,du
=\lambda(n_*)\Rz(n_*)
 \frac{\omega_{\Sset}(n_*)}{\Omega_{\Sset}(n_*)}.
\end{equation}
Each integral is $1/\Omega_{\Sset}(n_*)$, and $\lambda(p)=-1$.
The ratio is $1$ unless the square of an $\Sset$-prime divides $n_*$;
in that case the difference from the desired value has absolute value
at most $1$.
The resulting error after all bounded factors has squared norm
\begin{equation}\label{eq:square-error}
\ll YL^{C'}\sum_{p>\exp(L^{0.22})}p^{-2}
\ll YL^{C'}\exp(-L^{0.22}).
\end{equation}
For counts with floors,
$\sum_{p^2\le O(Y)}(Y/p^2+1)
 \ll Y\exp(-L^{0.22})+\sqrt Y$, with the latter term absorbed.
Remove this error globally by Parseval before estimating the arc
union. For the grid integral at fixed $Q$, arcs have spacing $1/Q$
and width $O(L^{C_0}/Y)$. Since $Y/Q=H'/r$ exceeds every fixed
power of $L$, they are disjoint for large $N$. The same norm bound
therefore handles this integral; summing in $r$ costs a fixed power.

After extraction the residual coefficient, for $w=n/p$, is
\begin{equation}\label{eq:residual-coefficient}
b_u(w)=\lambda(w_*)\Rz(w_*)u^{\Omega_{\Sset}(w_*)}
       \ind_{(w,N)=1}c(w_{\Pset}),\qquad |b_u(w)|\le L^C.
\end{equation}
Here $p\notin\Pset$, $p\nmid N$; no condition $p\nmid w$ is imposed.
Its discrepancy is exactly the square error already removed.

We need the principal residual mean. If $\log R\asymp L$, $\chi_0$
is principal modulo $Q=Nr$, and $I\subset[R,2R]$, then for every
fixed $B$ and $|t|\le L^B$,
\begin{equation}\label{eq:principal-residual}
\sum_{w\in I}b_u(w)\chi_0(w)w^{it}\ll_A RL^{-A}
\end{equation}
uniformly in $u\in[0,1]$ and $c$.
Write $w=gm$ with $g$ $\Pset$-smooth.
The function of $m$ is completely multiplicative: it is $0$ on
$\Pset$ and primes at most $z$, $-u$ on $\Sset$, and $-1$ elsewhere.
Lemma~\ref{lem:modifications} applies with $d=0.40$.
The mask $(m,N)=1$ is supplied by $\chi_0$.
For $g\le\sqrt R$ the residual length $R/g$ is a positive power of
$N$, and errors sum using
\begin{equation}\label{eq:smooth-harmonic}
\sum_{\substack{g\ {\Pset\text{-smooth}}}}\frac{|c(g)|}g
\le L^C\prod_{p\in\Pset}(1-1/p)^{-1}\le L^{C'}.
\end{equation}
For $g>\sqrt R$ use the trivial count and Rankin with
$\delta=L^{-0.40}$. Since all group primes are at most $\exp(L^{0.40})$,
the Euler product at $1-\delta$ is a fixed power of $L$.
The tail is $O(RL^{C'}\exp(-c_1L^{0.60}))$. If $\Pset$ is empty
the tail is zero. Differences of prefixes prove the assertion for
all interval endpoints. No multiplicativity of $c$ is needed.

Decompose $p$ and $w$ into dyadic intervals $p\asymp P$,
$w\asymp R$, with $PR\asymp Y$. There are $O(L^{C'})$ pieces and
\begin{equation}\label{eq:residual-length}
P\le\exp(L^{0.28}),\qquad
R/Q\gg H'/(Pr)\longrightarrow\infty
\end{equation}
faster than every fixed power of $L$. A factor $2$ in a prime dyad
endpoint is harmless by using a slightly wider fixed exponent range
in Lemma~\ref{lem:prime}.
For a separated power twist put
\begin{equation}\label{eq:character-polynomials}
P_\chi(t)=\sum_{p\in I_P\cap\Sset}\chi(p)p^{it},
\qquad R_\chi(t)=\sum_{w\in I_R}b_u(w)\chi(w)w^{it}.
\end{equation}
Orthogonality and the number of points in each residue class give,
for every real $t$,
\begin{equation}\label{eq:residual-energy}
\begin{split}
\frac1{\varphi(Q)}\sum_{\chi\bmod Q}|R_\chi(t)|^2
&=\sum_{a\in(\Z/Q\Z)^*}
  \left|\sum_{\substack{w\in I_R\\w\equiv a\ (Q)}}
                 b_u(w)w^{it}\right|^2\\
&\le(2R/Q+1)\sum_{w\in I_R}|b_u(w)|^2
 \ll R^2L^{C'}/Q.
\end{split}
\end{equation}
For $|t|$ bounded by any fixed power of $L$, Lemma~\ref{lem:prime}
and \eqref{eq:residual-energy} imply
\begin{equation}\label{eq:nonprincipal-energy}
\frac1{\varphi(Q)}\sum_{\chi\ne\chi_0}|P_\chi(t)R_\chi(t)|^2
\ll \frac{Y^2}{Q}P^{-1/8}L^{C'}.
\end{equation}
For $\chi_0$ use $|P_{\chi_0}|\le2P$,
\eqref{eq:principal-residual} and
$Q/\varphi(Q)\ll\log\log(3Q)$.
The full character energy is consequently
$O_A(Y^2L^{-A}/Q)$, without a triangle sum over $Q$ characters.

On a fixed product dyad, the factor
$\psi(pw/Y)(pw)^{iv}\e(\beta pw)$ is smooth in $\log(pw/Y)$.
For $|\beta|Y\le L^{C_0}$ its Mellin transform has integral norm
bounded by a fixed power of $L$.
Truncate at $|t|\le L^B$ with $B$ larger than the derivative exponents.
Repeated integration by parts gives any fixed logarithmic saving
in the tail. For every real $t$, the trivial $|P_\chi|\le2P$
and \eqref{eq:residual-energy} retain the baseline
$Y^2L^{C'}/Q$, so the tail loses no factor $Q$.
Minkowski in the normalized character Hilbert space, the $u$ integral
and the dyadic sum gives, for the extracted sequence $a_0$,
\begin{equation}\label{eq:full-character-energy}
\frac1{\varphi(Q)}\sum_{\chi\bmod Q}
 \left|\sum_na_0(n)\chi(n)\e(\beta n)\right|^2
\ll_A Y^2L^{-A}/Q.
\end{equation}
The supports are units modulo $Q$: every prime factor of $r$
is at most $L^{C_0}<z$, removed by roughness, and the $N$-unit mask
removes the others. Finite additive Parseval and character
orthogonality therefore give
\begin{equation}\label{eq:additive-character-parseval}
\begin{split}
\sum_{b\bmod Q}|\widehat a_0(b/Q+\beta)|^2
&=Q\sum_{a\bmod Q}
  \left|\sum_{n\equiv a\ (Q)}a_0(n)\e(\beta n)\right|^2\\
&=\frac Q{\varphi(Q)}
  \sum_{\chi\bmod Q}
  \left|\sum_na_0(n)\chi(n)\e(\beta n)\right|^2
 \ll_A Y^2L^{-A}.
\end{split}
\end{equation}
Integrate over length $2L^{C_0}/Y$, increasing the preceding precision,
and restore \eqref{eq:square-error}. This proves
\eqref{eq:grid-energy} for each fixed $N$.
\end{proof}

\begin{lemma}[Two-dimensional partial summation]\label{lem:abel}
The smooth-factor estimate \eqref{eq:full-character-energy} also
follows directly by Abel summation, with the same assumptions.
\end{lemma}
\begin{proof}
On a prime/residual rectangle $PR\asymp Y$ define
\begin{equation}\label{eq:abel-prefixes}
A_\chi(x)=\sum_{\substack{p\in I_P\cap\Sset\\p\le x}}\chi(p),
\qquad
B_\chi(y)=\sum_{\substack{w\in I_R\\w\le y}}b_u(w)\chi(w).
\end{equation}
Lemma~\ref{lem:prime} at $t=0$, \eqref{eq:principal-residual}
and \eqref{eq:residual-energy} hold uniformly in these prefixes.
Consequently
\begin{equation}\label{eq:prefix-energy}
\|(A_\chi(x)B_\chi(y))_\chi\|_{\ell^2(d\chi/\varphi(Q))}
\ll_A \frac{Y}{\sqrt Q}L^{-A}.
\end{equation}
The nonprincipal part has the stronger factor
$P^{-1/16}$ times a fixed power of $L$; the principal coordinate
uses any prescribed preceding precision.
Set
\begin{equation}\label{eq:abel-factor}
F(p,w)=\psi(pw/Y)(pw)^{iv}\e(\beta pw).
\end{equation}
Its two-variable variation satisfies
\begin{equation}\label{eq:abel-variation}
\begin{split}
|F(2P,2R)|&
+\int_P^{2P}|F_p(x,2R)|\,dx
+\int_R^{2R}|F_w(2P,y)|\,dy\\
&+\int_P^{2P}\int_R^{2R}|F_{pw}(x,y)|\,dy\,dx
\ll L^{C'}.
\end{split}
\end{equation}
Indeed rescale $p=Ps,w=Rt$. On $[1,2]^2$, $PR/Y$ is in a fixed
compact subset of $(0,\infty)$; derivatives through order two are
bounded by a fixed polynomial in the derivative norms of $\psi$,
$|v|$ and $|\beta|Y$. The factor $(PR)^{iv}$ has modulus $1$.

With prefixes zero at the lower endpoints, two-dimensional partial
summation is the exact identity
\begin{equation}\label{eq:abel-identity}
\begin{split}
S_\chi={}&F(2P,2R)A_\chi(2P)B_\chi(2R)\\
&-\int_P^{2P}F_p(x,2R)A_\chi(x)B_\chi(2R)\,dx\\
&-\int_R^{2R}F_w(2P,y)A_\chi(2P)B_\chi(y)\,dy\\
&+\int_P^{2P}\int_R^{2R}
       F_{pw}(x,y)A_\chi(x)B_\chi(y)\,dy\,dx,
\end{split}
\end{equation}
where $S_\chi$ is the weighted double character sum.
Half-open dyads can be handled by displaced real endpoints.
Minkowski, \eqref{eq:prefix-energy} and \eqref{eq:abel-variation}
retain $Y/\sqrt Q$ while costing only a fixed power of $L$.
Integrating $u$ and summing the polynomially many dyads proves
\eqref{eq:full-character-energy} uniformly in $\beta$.
\end{proof}

\begin{remark}\label{rem:raw-scale}
The condition $R=Y/P\gg Q$ is essential. At $H'=1$, the
orthogonality bound contains $1+Q/R$; with $Q\asymp N$ and $Y=N$
the nonprincipal saving becomes
$P^{-1/8}(1+Q/R)\asymp P^{7/8}$.
The grid lemma is used only for comparisons with two nonempty
free products, never for the raw correlation.
The narrow displacement $|\beta|Y\le L^{C_0}$ is precisely what
makes \eqref{eq:abel-variation} polynomial in $L$.
\end{remark}

\section{Reflected weighted cancellation}\label{sec:cancellation}

\begin{definition}[Marked weights]\label{def:weights}
For nonnegative integers $\mathbf t=(t_g)_g$ let $(y)_t$
denote the falling factorial and put
\begin{equation}\label{eq:marked-weight}
\wt(n)=q_0^{\omega_{\Pset}(n)-\sum_gt_g}
            \prod_g\frac{(\omega_g(n))_{t_g}}{V_g^{t_g}}.
\end{equation}
More generally, replace the product of falling factorials by the sum
of a mark function $b$ over the ordered distinct unshared divisor
lists, retaining the same normalization and damping; write
$W_{\mathbf t}^{b}(n)$.
An empty admissible list set gives weight zero. If $\max t_g\le t_{\max}$
is fixed, then
\begin{equation}\label{eq:marked-bound}
\sup_n|W_{\mathbf t}^{b}(n)|\le\|b\|_\infty L^{C(t_{\max})}.
\end{equation}
Indeed $\sup_{y\ge t}q_0^{y-t}(y)_t/V_g^t$ is a fixed constant
for $t\le t_{\max}$, and there are $O(\log L)$ groups.
For $\ell=1$ write
\begin{equation}\label{eq:raw-weight}
W_1(n)=q_0^{\omega_{\Pset}(n)-s}\prod_g\frac{\omega_g(n)}{V_g},
\qquad
H(n)=\Rz(n)\ind_{(n,N)=1}\ind_{\omega_{\Sset}(n)>0}W_1(n).
\end{equation}
\end{definition}

\begin{proposition}[Reflected cancellation]\label{prop:cancellation}
Use the ranges \eqref{eq:separated-ranges}, with $O(\log L)$ disjoint
groups, reciprocal masses in a fixed compact subset of $(0,\infty)$,
and at least $c_0\log L$ groups in each of the small and big ranges.
No group prime divides $N$. Assume every big group is a prime
interval with the primes dividing $N$ removed. Fix $0<q_0<1$ and
$\ell=1$. For every fixed smooth $\psi_0$ supported in $[1/3,2/3]$
and every fixed $A>0$,
\begin{equation}\label{eq:reflected-cancellation}
\left|\sum_{1\le n<N}\psi_0(n/N)H(n)H(N-n)
                       \lambda(n)\lambda(N-n)\right|
\ll_A NL^{-A}.
\end{equation}
The threshold is uniform over $N$ and the indicated group construction; no average over $N$ is taken.
For odd $N$ the rough support is empty.
\end{proposition}
\begin{proof}
Set $x=N$, take $k=1$, and choose the residual family of
Input~\ref{input:ik}, using Proposition~\ref{prop:punctures}.
Use $B_N$ from Proposition~\ref{prop:reflection}.
We verify the lift, the individual patterns, and then their comparison
terms.

At a positive physical position use, at the source and at the target,
the mark-independent endpoint function
\begin{equation}\label{eq:physical-endpoint}
f(n)=\lambda(n)F_0(n)q_0^{(\omega_{\Pset}(n)-Ms)/2},
\qquad
F_0(n)=\Rz(n_*)\ind_{(n,N)=1}
                         \ind_{\omega_{\Sset}(n_*)>0}.
\end{equation}
Both endpoints are zero outside their positive cutoffs.
Every group prime exceeds $z$, is outside $\Sset$, and is coprime
to $N$, so $F_0$ is invariant under multiplication by
$\Pset$-integers. On a physical state $\omega_{\Pset}(n)\ge Ms$,
and therefore $|f(n)|\le1$.

For a positive dyad, summing first over ordered lists, then over
multiples of their product $P_{\rm list}$, gives
\begin{equation}\label{eq:physical-norm}
\begin{split}
\sum_{n\asymp Y}\sum_{\text{lists at }n}\prod_gV_g^{-M}
&\le CY\prod_g
  \left(\sum_{p\in\Pset_g}\frac1{pV_g}\right)^M+N^{o(1)}
\ll Y,\\
\log P_{\rm list}&=O_J(L^{0.40}\log L)=o(L).
\end{split}
\end{equation}
The rounding term is $N^{o(1)}$, since $J$ is fixed.
Thus both endpoint norms are $O(\sqrt Y)$, with the logarithmic
exponent independent of $J$ and the probe count.
Bounded smooth cutoffs and modulus-one power twists preserve this.

Insert the edge factor $\Psi(n/(ND))/n$, where
\begin{equation}\label{eq:lift-cutoff}
\Psi(t)=t\psi_0(t).
\end{equation}
Its support makes $n'=ND-n$ positive and comparable to $n$.
Since $D=\exp(O_J(L^{0.40}\log L))$, only a fixed power of $L$
many size dyads are needed, and all their lengths lie between
$N^\gamma$ and $N^{1/\gamma}$ for a fixed $\gamma>0$.

Choose a smooth dyadic partition of unity and smooth target cutoffs
equal to $1$ on the corresponding comparable positive target ranges.
If the source scale is $Y$, the source and target cutoff derivatives
of order $j$ are $O_j(Y^{-j}L^{C(j+1)})$.
They may be obtained by rescaling fixed smooth functions; the
logarithmic derivative exponents are fixed before the graph parameters.
Overlapping size intervals may each be split into finitely many
support dyads for transference. The smooth partition and target
cutoffs sum back to the original expression on the support of $\Psi$
before any individual pattern is evaluated. All subsequent cutoffs
used for Lemmas~\ref{lem:grid} and~\ref{lem:abel} have these same
derivative bounds.

For $\Phi(u)=\Psi(\exp u)$ and
$\widehat\Phi(t)=\int_\R\Phi(u)\exp(-itu)\,du$, logarithmic Fourier inversion is
\begin{equation}\label{eq:lift-mellin}
\Psi(n/(ND))=\frac1{2\pi}\int_\R
 \widehat\Phi(t)(n/N)^{it}D^{-it}\,dt.
\end{equation}
On a source dyad place $Y/n$ and the conjugated source twist in
the first endpoint, and divide the pairing by $Y$.
The allowed ideal frequency is $\zeta=-t$.
The transform has bounded or fixed-log-power integral norm.
If its logarithmic derivative exponent is $C_\psi$, fix
$F>C_\psi+2$ in advance and take $C_4=F$ in Input~\ref{input:ik}, so
that every frequency $\zeta=-t$ with $|t|\le L^F$ is admissible there.
Integration by parts gives
\begin{equation}\label{eq:lift-tail}
\int_{|t|>L^F}|\widehat\Phi(t)|\,dt
\ll_j L^{C_\psi(j+1)-F(j-1)}.
\end{equation}
Arbitrarily strong saving follows by increasing $j$ without changing
$F$ or $C_4$. By Propositions~\ref{prop:dilation}
and~\ref{prop:reflection}, the absolute row and column sums of $B_N$ are
at most $L^{C_{\rm abs}}$ for every $\zeta$, so each pairing is
$O(YL^{C_{\rm abs}})$ by the endpoint norms; together with
\eqref{eq:lift-tail} this disposes of the tails.
Choose the residual precision, then $E_{\rm id}$, then the probes and
kernel precisions, then sufficiently large fixed $J$, and finally $N$.
Propositions~\ref{prop:dilation}--\ref{prop:punctures} and
Input~\ref{input:ik} make the entire normalized signed residual
pairing $O(L^{-B})$ for any prescribed $B$.

We now compute its individual patterns. This uses the normalization
of [S, Section~5.2, pp.~36--38] and the shared-sign identity of
[E, (4.18)--(4.19), p.~16], with reflection included explicitly.
In group $g$ put $c_g=|C_g|$, $s_{R,g}=J-c_g$,
$t_g=1+c_g$, so $M=s_{R,g}+t_g$.
Write $D=D_RD_C$ for shared and auxiliary free products.
Shared labels divide both endpoints; hence
\begin{equation}\label{eq:shared-reflection}
n=D_Rw,\qquad
n'=ND-n=D_Rw',\qquad
w+w'=ND_C.
\end{equation}
The cutoff and harmonic normalization become exactly
\begin{equation}\label{eq:shared-cutoff}
\frac{\Psi(n/(ND))}{n}
=\frac{\psi_0(w/(ND_C))}{ND_RD_C}.
\end{equation}
The state and target-list normalizations satisfy
\begin{equation}\label{eq:shared-normalization}
V_g^{-M-t_g}=V_g^{-s_{R,g}-2t_g}
=V_g^{-s_{R,g}}V_g^{-t_g}V_g^{-t_g}.
\end{equation}
The factor $1/D_R$ converts every shared-label sum with normalization
$V_g^{-1}$ into a probability draw $1/(pV_g)=\mu_g(p)$, leaving
$V_g^{-t_g}$ at each endpoint.

Away from shared overlaps,
\begin{equation}\label{eq:shared-degrees}
\omega_g(D_Rw)-M=\omega_g(w)-t_g,\qquad
\omega_g(D_Rw')-M=\omega_g(w')-t_g.
\end{equation}
At each endpoint the graph half-damping in
\eqref{eq:physical-multiplier} and the endpoint half-damping in
\eqref{eq:physical-endpoint} give
$q_0^{\omega_{\Pset}(n)-Ms}$.
By \eqref{eq:shared-degrees} these are exactly the full powers
$q_0^{\omega_{\Pset}(w)-\sum_gt_g}$ and
$q_0^{\omega_{\Pset}(w')-\sum_gt_g}$ in the two marked weights.
The sign identity, valid even on overlaps, is
\begin{equation}\label{eq:shared-sign}
\lambda(D_Rw)\lambda(D_Rw')
=\lambda(D_R)^2\lambda(w)\lambda(w')
=\lambda(w)\lambda(w').
\end{equation}
Also $F_0(D_Rw)=F_0(w)$ and $F_0(D_Rw')=F_0(w')$.
Comparison coefficients use free and designated unshared big labels,
never shared labels, so they are unchanged by summing the shared draws.
The endpoint vectors are mark-independent, so symmetrization is an
average that does not alter this calculation.

We justify replacement by unrestricted shared draws in absolute
value. Given $w,w',D_C$ and preceding shared draws, at most
$O_J(L)$ primes are forbidden: the prime divisors of the two endpoints,
free labels and preceding shared labels.
Every draw has atom $O(\exp(-L^{0.12}))$; there are $O_J(\log L)$
draws. For distinct shared values, the exact overlap identity is
\begin{equation}\label{eq:overlap-degree}
\omega_{\Pset}(D_Rw)-Ms
=\omega_{\Pset}(w)-\sum_gt_g
-\#\{p:p\mid D_R,\ p\mid w\}.
\end{equation}
Ignoring the overlap damping loses at most
$q_0^{-2Ms}=L^{O_J(1)}$ across the two endpoints.
Unshared divisor lists still divide $w,w'$.
Their bounded marked weights in \eqref{eq:marked-bound}, $|F_0|\le1$,
and harmonic normalization bound the remaining absolute sum by a
fixed power of $L$.
The total shared-exclusion error is
\begin{equation}\label{eq:shared-error}
L^{O_J(1)}\exp(-L^{0.12}).
\end{equation}
If a banned free prime $p$ divides an endpoint, then $p\mid D_C$
and \eqref{eq:shared-reflection} imply it divides both $w,w'$.
At scale $Y\asymp ND_C$, summing these multiples costs
$O(YL^C/p)$, or $O(L^C/p)$ after normalization.
Indeed $w=pv$, $w'=p(ND_C/p-v)$; both ranges are positive and
have length $O(Y/p)$, with $Y/p=N^{1+o(1)}$.
There are only $O(\log L)$ free slots with fixed probe count and
$p\ge\exp(L^{0.35})$, so this error is also smaller than every
fixed negative power of $L$. No cancellation is used for deleted terms.

After these replacements the pattern is
\begin{equation}\label{eq:factored-pattern}
\frac1N\E_{\rm free}\frac1{D_C}
\sum_{\substack{w>0\\w'=ND_C-w>0}}
 \psi_0(w/(ND_C))F_0(w)F_0(w')
 \lambda(w)\lambda(w')\wt(w)\wt(w')
 \E_{\rm marks}K_\nu .
\end{equation}
The inner expectation averages independent uniform ordered unshared
divisor lists, giving zero if either list set is empty.
The raw pattern has $D_C=1$, $t_g=1$, $K_0=1$, and gives exactly
\begin{equation}\label{eq:primitive-pattern}
\frac1N\sum_{1\le w<N}\psi_0(w/N)
H(w)H(N-w)\lambda(w)\lambda(N-w).
\end{equation}

For each comparison expand the two cell/character kernels in
\eqref{eq:cell-kernel} and separate the smooth factors logarithmically.
Up to a total fixed power of $L$ in coefficients and integral norms,
a term indexed by a pair of cells $(j,j')$ is
\begin{equation}\label{eq:comparison-term}
\frac1Y\E_{D_{\mathcal A},D_{\mathcal B}}
 b_{\mathcal A}(D_{\mathcal A})b_{\mathcal B}(D_{\mathcal B})
 \sum_w U(w)V(ND_{\mathcal A}D_{\mathcal B}-w),
\quad |b_{\mathcal A}|,|b_{\mathcal B}|\le1.
\end{equation}
For this term let $U_{\mathcal A},U_{\mathcal B}$ be the lower
multiplicative size endpoints of the two cells, and put
$H'=U_{\mathcal A}U_{\mathcal B}$, $Y=NH'$.
Explicitly, the term belongs to primitive characters $\chi,\chi'$ of
conductor at most $Q_0$ and to a real frequency $\tau$, $|\tau|\le L^F$,
of the expansion \eqref{eq:lift-mellin} applied to $\psi_0$ in place of
$\Psi$, and
\begin{equation}\label{eq:comparison-data}
\begin{gathered}
b_{\mathcal A}(D)=\ind_{\log D\in I_j}\overline{\chi(D)}\,
 \frac{U_{\mathcal A}}{D}\,D^{-i\tau},\qquad
b_{\mathcal B}(D)=\ind_{\log D\in I_{j'}}\overline{\chi'(D)}\,
 \frac{U_{\mathcal B}}{D}\,D^{-i\tau},\\
U(w)=\Xi(w/Y)(w/N)^{i\tau}\lambda(w)F_0(w)W_{\mathbf t}^{b_{\rm s}}(w),\qquad
V(w')=\Xi(w'/Y)\lambda(w')F_0(w')W_{\mathbf t}^{b_{\rm t}}(w').
\end{gathered}
\end{equation}
Here $\Xi$ is a fixed smooth function with compact support in
$(0,\infty)$ and $\Xi=1$ on $[1/3,2\mathrm e^2/3]$, the source mark
function is
$b_{\rm s}=\ind_{\log X_{\mathcal B}\in I_{j'}}\chi'(X_{\mathcal B})$, and
the target mark function is
$b_{\rm t}=\ind_{\log X_{\mathcal A}\in I_j}\chi(X_{\mathcal A})$.
On the two cells $D_{\mathcal A}D_{\mathcal B}\in[H',\mathrm e^2H')$, since
$\mathfrak h\le1$. As $\psi_0$ is supported in $[1/3,2/3]$, on the support of
$\psi_0(w/(ND_{\mathcal A}D_{\mathcal B}))$ both $w$ and
$w'=ND_{\mathcal A}D_{\mathcal B}-w$ lie in $[Y/3,2\mathrm e^2Y/3]$, so both
cutoffs equal $1$ there. Moreover
$1/(ND_{\mathcal A}D_{\mathcal B})
 =Y^{-1}(U_{\mathcal A}/D_{\mathcal A})(U_{\mathcal B}/D_{\mathcal B})$.
The term carries the coefficient
$-(-1)^{|\mathcal A|+|\mathcal B|}\widehat{\Phi_0}(\tau)/(2\pi\nu_j\nu_{j'})$,
where $\Phi_0(u)=\psi_0(\exp u)$ and $\nu_j,\nu_{j'}\ge\xi$ are the two
cell probabilities; the frequencies $|\tau|>L^F$ are discarded as in
\eqref{eq:lift-tail}, using $\|U\|_2\|V\|_2\ll YL^C$.
The free products are independent, nonempty, and each contains a big
prime. Thus
\begin{equation}\label{eq:comparison-scales}
\exp(2L^{0.35})/4\le H'\le
             \exp(C_2L^{0.40}\log L),\qquad
\rho_{\mathcal A}(d),\rho_{\mathcal B}(d)\ll L^C/d.
\end{equation}
The harmless factor $1/4$ covers endpoint conventions.
The mass bounds are \eqref{eq:product-collapse}, with at most the
fixed probe count per group.
Both endpoint sequences have squared norm $O(YL^C)$.
They have exactly the form \eqref{eq:grid-sequence}: the group
coefficient is $\lambda(g)$ times its bounded marked weight and
big-mark kernel factors. These depend only on $g=n_{\Pset}$.
The first endpoint retains a reserved $\Sset$-prime.

Handle the major arcs separately for this particular $(j,j')$ and
these $U,V,Y$. Fourier inversion supplies
\begin{equation}\label{eq:comparison-multiplier}
\mathcal M_{j,j'}(\theta)=
\E b_{\mathcal A}(D_{\mathcal A})b_{\mathcal B}(D_{\mathcal B})
       \e(-N\theta D_{\mathcal A}D_{\mathcal B}).
\end{equation}
Apply Dirichlet approximation to $N\theta$ with maximum denominator
$\lfloor H'/L^B\rfloor$.
For its reduced fraction $u/r$ the error is
$O(L^B/(rH'))$ and at most $r^{-2}$.
If $r>L^B$, the bilinear estimate [S, Lemma~4.2, (4.5)--(4.6),
pp.~28--29] and \eqref{eq:comparison-scales} give arbitrary fixed
logarithmic saving. More explicitly, its squared normalized bound
is a fixed power of $L$ times
\[
\frac1{U_{\mathcal B}}+\frac1r
+\frac{\log(2r)}{U_{\mathcal A}}
+\frac{r\log(2r)}{H'}.
\]
Both product sizes exceed $\exp(L^{0.35})/2$; the log factor is a
fixed power of $L$, and $r\le H'/L^B$. Increasing $B$ therefore gives
$|\mathcal M_{j,j'}|\ll L^{-D}$ for any prescribed $D$.
For $r\le L^B$ one has
\begin{equation}\label{eq:comparison-arcs}
\theta=\frac b{Nr}
 +O\left(\frac{L^B}{NrH'}\right),
\end{equation}
so the major set $\mathfrak M_{j,j'}$ is contained in the grids of
Lemma~\ref{lem:grid} with width at most $L^B/Y$.
Its many rational centers are retained in that lemma.

On the complement of $\mathfrak M_{j,j'}$, Parseval bounds
\eqref{eq:comparison-term} by $L^{-D+C}$.
On $\mathfrak M_{j,j'}$ use $|\mathcal M_{j,j'}|\le1$,
Lemma~\ref{lem:grid} for $U$, and global Parseval for $V$:
\begin{equation}\label{eq:major-cauchy}
\frac1Y\int_{\mathfrak M_{j,j'}}
          |\widehat U(\theta)\widehat V(\theta)|\,d\theta
\ll_A L^{-A/2+C}.
\end{equation}
Reflection only changes the Fourier conjugation convention, not this
absolute bound. Choose the grid precision after all fixed comparison
costs, and sum these estimates term by term over the cell pairs,
characters and separated smooth integrals.
Subtract the comparisons from the small residual pairing and restore
the two exclusion errors and Fourier tails.
The raw expression \eqref{eq:primitive-pattern} proves
\eqref{eq:reflected-cancellation}.

Finally, $H$ depends only on the fixed groups, $q_0$ and $\ell=1$,
not on $J$. Increasing accuracy changes the auxiliary graph and
threshold but leaves this weight fixed.
The comparison coefficient exponents, set by the probes, are paid
afterward by the arbitrary-precision grid lemma.
No choice of $J$ feeds back into the ideal norm target.
\end{proof}

\section{A bounded weight with positive rough-pair mass}\label{sec:weight}

\begin{proposition}[Weight construction and moments]\label{prop:mass}
In each range
$[\exp(L^{0.12}),\exp(L^{0.18})]$ and
$[\exp(L^{0.35}),\exp(L^{0.40})]$, list the primes not dividing $N$
in increasing order and close a group when its reciprocal mass reaches
$10$. Discard the last incomplete group in each range. Then
\begin{equation}\label{eq:constructed-groups}
\begin{gathered}
10\le V_g\le10+\exp(-L^{0.12}),\\
|\mathcal G_{\rm small}|=0.006\log L+O(1),\qquad
|\mathcal G_{\rm big}|=0.005\log L+O(1).
\end{gathered}
\end{equation}
Use $q_0=9/10$, $\ell=1$ and \eqref{eq:raw-weight}.
Fix a smooth $0\le\psi_0\le1$ supported inside $[1/3,2/3]$ with
$I_\psi=\int\psi_0(t)\,dt>0$.
For sufficiently large even $N$,
\begin{equation}\label{eq:positive-mass}
\sum_{1\le n<N}\psi_0(n/N)H(n)H(N-n)
\ge cNL^{-59/250},\qquad 59/250=0.236.
\end{equation}
For every fixed $A>0$ the two one-sign moments satisfy
\begin{equation}\label{eq:one-sign-moments}
\begin{split}
\sum_{1\le n<N}\psi_0(n/N)H(n)H(N-n)\lambda(n)&=O_A(NL^{-A}),\\
\sum_{1\le n<N}\psi_0(n/N)H(n)H(N-n)\lambda(N-n)&=O_A(NL^{-A}).
\end{split}
\end{equation}
Moreover $0\le H(n)\le1$ for every $n$.
\end{proposition}
\begin{proof}
Mertens' harmonic prime estimate and \eqref{eq:deleted-mass} give
\eqref{eq:constructed-groups} uniformly in $N$.
The overshoot at closing a group is at most the largest reciprocal
prime in its range, and the total overshoot is $o(1)$.
For large $N$, $s\le0.012\log L$ and each family contains at least
$0.004\log L$ groups.
Thus all graph inputs hold with the fixed choices
\begin{equation}\label{eq:numerical-graph-data}
\begin{gathered}
a=0.12,\quad b=0.18,\quad c=0.35,\quad d=0.40,\\
c_0=0.004,\quad C_0=0.02,\quad v_-=10,\quad v_+=11,\quad
q_0=0.9,\quad\ell=1.
\end{gathered}
\end{equation}
Every big group is an interval with the primes dividing $N$ removed.
All these data are fixed before $N$; only allowed prime sets,
group counts and the source moment length grow.

For $k\ge1$,
\begin{equation}\label{eq:pointwise-local}
kq_0^{k-1}\le1+q_0+\cdots+q_0^{k-1}
\le(1-q_0)^{-1}=10\le V_g.
\end{equation}
Each local factor in $W_1$ is at most $1$, and a missing group
gives $0$. The remaining factors of $H$ are indicators.

Temporarily omit the $N$-unit masks and use independent uniform
residues at each prime. The rough factor is
\begin{equation}\label{eq:rough-euler}
V_N(z)=\prod_{p\le z}(1-\nu_N(p)/p),\qquad
\nu_N(p)=
\begin{cases}1,&p\mid N,\\2,&p\nmid N.\end{cases}
\end{equation}
For even $N$ the factor at $2$ is $1/2$. For $p\ge3$, use
$1-2/p$ as a lower bound. Its ratio to $(1-1/p)^2$ has a positive
convergent product, so
$V_N(z)\ge c/(\log z)^2=cL^{-1/5}$ uniformly in $N$.

At a group prime $p\nmid N$, the two endpoint divisibility indicators
cannot both be $1$. The degree generating function is
\begin{equation}\label{eq:group-generating}
J_g(u,v)=\prod_{p\in\Pset_g}
            \left(1+\frac{u-1}{p}+\frac{v-1}{p}\right).
\end{equation}
The local expectation of the two $W_1$ factors is exactly
\begin{equation}\label{eq:group-local-mass}
\begin{split}
h_g&=V_g^{-2}\partial_u\partial_vJ_g(q_0,q_0)\\
&=\frac{J_g(q_0,q_0)}{V_g^2}
 \left\{\left(\sum_{p\in\Pset_g}\frac1{p+2q_0-2}\right)^2
       -\sum_{p\in\Pset_g}\frac1{(p+2q_0-2)^2}\right\}.
\end{split}
\end{equation}
Since $V_g=10+o(1)$ uniformly and the smallest group prime tends
to infinity,
\begin{equation}\label{eq:group-mass-asymptotic}
h_g=\exp(2(q_0-1)V_g)(1+o(1))
=\mathrm e^{-2}(1+o(1))\ge\mathrm e^{-3}
\end{equation}
for all groups at sufficiently large $N$.
The error is $O(\sum_{p\in\Pset_g}p^{-2})$, with constants uniform
for $V_g\in[10,11]$.

The reserved-prime condition at both endpoints has factor
\begin{equation}\label{eq:reserved-mass}
B_{\Sset}=1-2\prod_{p\in\Sset}(1-1/p)+\prod_{p\in\Sset}(1-2/p)
\ge1-2\prod_{p\in\Sset}(1-1/p)\ge\frac12.
\end{equation}
Here $\sum_{p\in\Sset}1/p=0.06\log L+O(1)$.
The rough primes, graph groups and reserved primes are disjoint.
Their independent mean is consequently
\begin{equation}\label{eq:independent-mass}
\mathfrak M_N=V_N(z)B_{\Sset}\prod_gh_g
\ge cL^{-1/5}\mathrm e^{-3s}
\ge cL^{-1/5-0.036}=cL^{-59/250}.
\end{equation}

We now justify a finite residue comparison strong enough to multiply
by signs. Each endpoint group factor expands exactly as
\begin{equation}\label{eq:local-residue-expansion}
\frac1{V_g}\sum_{p\in\Pset_g}\ind_{p\mid n}
 \prod_{\substack{p'\in\Pset_g\\p'\ne p}}
       (1-(1-q_0)\ind_{p'\mid n}),
\end{equation}
and likewise for $N-n$. The selected mark primes at the two
endpoints must be distinct because $p\nmid N$.
There are $2s$ selected primes in every nonzero term, with product
at most $\exp(O(L^{0.40}\log L))$.
The sum of their normalized independent congruence densities is at
most $\prod_g(V_g/V_g)^2=1$.

Write each reserved-prime indicator as $1$ minus its absence product
and expand into four terms. Conditional on any selected mark tuple,
each resulting unsigned avoidance product has factors $1-X_p$,
where $0\le X_p\le1$, and independent expected sum
$\sum_p\E X_p=O(\log L)$ uniformly in that tuple.
At group primes the two endpoint predicates are disjoint.
At a rough prime dividing $N$ the coincident predicates are counted
only once. Evaluate selected-prime factors before expanding.
Bonferroni truncations at consecutive orders $h,h+1$, with $h$ an
even integer exceeding $C_A\log L$, sandwich these products.
The independent expected gap is at most
\begin{equation}\label{eq:bonferroni-gap}
\frac{(C\log L)^{h+1}}{(h+1)!}\le L^{-A}
\end{equation}
after $C_A$ is fixed sufficiently large.

Every retained term is a residue class, or a union expanded into at
most $2^{h+1}$ residue classes, modulo a squarefree product of at most
$2s+h+1$ primes at most $\exp(L^{0.40})$.
The total absolute coefficient mass, including selected marks and
residue choices, is
\begin{equation}\label{eq:residue-cost}
\exp(O_A(L^{0.40}\log L))=N^{o(1)}.
\end{equation}
For each residue class modulo $d$, smooth counting gives
$NI_\psi/d+O(1)$, uniformly in the residue and $d$.
One way to see the $O(1)$ is to integrate the bounded discrepancy
of a progression against the fixed function $\psi_0$.
The finite polynomial average therefore equals its independent
average with error $N^{-1+o(1)}$.
Combining the sandwich bounds and \eqref{eq:bonferroni-gap}
gives average error $O_A(L^{-A})+N^{-1+o(1)}$.
More strongly, choose one sandwich polynomial for each of the four
unsigned pieces. The nonnegative gap controls its absolute error.
Hence the exact nonnegative weighted product, without the $N$-unit
masks, has a finite residue polynomial approximation whose $L^1$
error against $\psi_0(n/N)$ is $O_A(NL^{-A})$.
The four pieces cost only a factor four.

Restoring the $N$-unit masks changes the bounded weighted sum only
if a prime $p\mid N$, $p>z$, divides an endpoint.
Its total contribution is at most
\begin{equation}\label{eq:unit-mask-error}
2\sum_{\substack{p\mid N\\p>z}}(N/p+1)
\ll \frac{NL}{z\log z}+\frac L{\log z},
\end{equation}
smaller than every $NL^{-A}$.
Thus the unsigned moment is
$NI_\psi\mathfrak M_N+O_A(NL^{-A})$,
which proves \eqref{eq:positive-mass}.

For a one-sign moment, remove the unit masks from the exact bounded
weight using \eqref{eq:unit-mask-error}, before passing to the residue
polynomial. Multiply its $L^1$ approximation by $\lambda(n)$.
All retained moduli are $N^{o(1)}$.
For a class $n\equiv a\pmod d$, put $g=(a,d)$.
Writing $n=gu$ leaves a reduced class modulo $d/g$ and
$\lambda(gu)=\lambda(g)\lambda(u)$.
Character orthogonality in the reduced class and
Lemma~\ref{lem:power}, with a fixed margin and partial summation for
$\psi_0$, give $O(N^{15/16+o(1)})$ per retained class.
The total coefficient mass \eqref{eq:residue-cost} makes the
aggregate $O(N^{0.99})$ for sufficiently large $N$.
Add the approximation error at arbitrary fixed precision.
The second one-sign moment follows by reflection and the cutoff
$\psi_0(1-t)$. This proves \eqref{eq:one-sign-moments}.
\end{proof}

\section{Proof of the two theorems}\label{sec:theorems}
\begin{proof}[Proof of Theorem~\ref{thm:main}]
Take $H,\psi_0$ from Proposition~\ref{prop:mass} and set
\begin{equation}\label{eq:signed-weight-count}
T_{\sigma,\tau}=\sum_{1\le n<N}\psi_0(n/N)H(n)H(N-n)
 \frac{1+\sigma\lambda(n)}2\frac{1+\tau\lambda(N-n)}2.
\end{equation}
Each summand is nonnegative, at most $1$, and supported on the set
in \eqref{eq:main}. Expanding the sign indicators gives one quarter
of the unsigned moment and the two one-sign moments in
Proposition~\ref{prop:mass}, and one quarter of the double-sign
moment in Proposition~\ref{prop:cancellation}.
For example, use precision $A=2$. Uniformly in all four signs,
\begin{equation}\label{eq:signed-lower}
T_{\sigma,\tau}\ge c_1NL^{-59/250}-C_1NL^{-2}
\ge c_2NL^{-59/250}.
\end{equation}
Since $59/250<1/4$, the unweighted count is at least the bound in
\eqref{eq:main}. All constants and the threshold are uniform in
even $N$ and the sign pair.
\end{proof}

\begin{remark}\label{rem:parity}
The parity information uses the graph residual estimate, complete
multiplicativity in \eqref{eq:shared-sign}, and the prime and
all-character estimates behind the grid lemma.
Unsigned residue data alone do not bound the double-sign moment.
Replacing every sign by $+1$ would fail
\eqref{eq:principal-residual}.
For a completely multiplicative sign with
$f(p)=(p/q)$ for $p\ne q$, $f(q)=-1$, and $q\equiv3\pmod4$ prime,
where $(\cdot/q)$ is the Legendre symbol,
a unit pair at an even total divisible by $q$ satisfies
\begin{equation}\label{eq:character-model}
f(n)f(N-n)=(n/q)(-n/q)=-1.
\end{equation} The shared-factor identity still holds, but the
reserved prime values are not constantly $-1$ and twisting by the
character modulo $q$ violates \eqref{eq:modified-mean}.
Thus the general sign function fails at the prime extraction and
principal residual step. The all-modulus zero-free input also
excludes any real zero above $7/8$ that would obstruct the character
contours. No character at the grid modulus is suppressed.
\end{remark}

Roughness alone gives, for either represented $a\in\{n,N-n\}$,
\begin{equation}\label{eq:rough-omega}
\Omega(a)\log z<\log a<\log N,\qquad
\Omega(a)<(\log N)^{9/10}.
\end{equation}
Thus $\lceil(\log N)^{9/10}\rceil$ already bounds both factor
counts with each prescribed parity. The next argument gives the
smaller growing bound \eqref{eq:factor-cutoff}.

\begin{proof}[Proof of Theorem~\ref{thm:postselection}]
The number of $z$-rough nonsquarefree integers $a\le N$ is at most
\begin{equation}\label{eq:nonsquarefree-count}
\sum_{p>z}\left\lfloor\frac N{p^2}\right\rfloor
\le N\sum_{m>z}m^{-2}\le\frac N{z-1}.
\end{equation}
For squarefree $a$, the divisor function satisfies
$d(a)=2^{\Omega(a)}$. The elementary divisor identity gives
\begin{equation}\label{eq:divisor-sum}
\sum_{a\le N}d(a)=\sum_{m\le N}\lfloor N/m\rfloor
\le N(1+\log N).
\end{equation}
The number of squarefree $a\le N$ with $\Omega(a)>K(N)$ is at most
\begin{equation}\label{eq:omega-tail}
2^{-K(N)}N(1+\log N)\ll NL^{1-2\log2}.
\end{equation}
Apply both bounds to each endpoint via $a\mapsto N-a$.
The total number of discarded representations, irrespective of sign,
is
\begin{equation}\label{eq:postselection-error}
\ll N/(z-1)+NL^{1-2\log2}=o(NL^{-1/4}),
\end{equation}
since $\log2>2/3$ and $1-2\log2<-1/3$.
Subtract this unsigned error separately from each of the four lower
bounds in Theorem~\ref{thm:main}. At least half its stated mass remains
for large $N$, giving the asserted possibly smaller constant.
\end{proof}

\bibliographystyle{plainnat}
\bibliography{references}
\section*{Use of AI}
The proofs and the text of this paper were produced with AI systems (GPT-6 Astra, GPT-6.1 Sol and Claude
Opus~5.5) under the author's direction. They were checked by independent reviews with AI models; the
zero-free input is a Lean theorem of the OpenAI library. The author takes responsibility for the
content.
\end{document}
